%!TEX root = ./../../thesis-trier.tex

The Rystsov graph and $\Gamma_1(\mathcal A)$
were introduced in~\cite{Rystsov2000} and~\cite{BondarV16}, respectively.
The graphs $\Gamma_k(\mathcal A)$
for $k \ge 1$ are from~\cite{BondarV18,BondarCV23}. The graph $\Gamma_1(\mathcal A)$
was referred to as the Rystsov graph in~\cite{DBLP:conf/latin/CasasV22}.  Here, however,
we reserve the term Rystsov graph for the graph $I(\mathcal A)$ originally 
introduced by Rystsov in~\cite{Rystsov2000}. 
In~\cite{DBLP:journals/ijfcs/GonzeGJGV19} an approach utilizing 
a sequence of increasing subgraphs of $\Gamma_1(\mathcal A)$
was used to prove a quadratic upper bound
on the reset length of automata whose transformation semigroup is the full 
transformation semigroup.

Gonze and Jungers~\cite{DBLP:journals/jalc/GonzeJ19} gave an algorithm running in time $O(|\Sigma||Q|^2)$
to compute $\Gamma_1(\mathcal A)$. They also refuted the conjecture of Bondar and Volkov~\cite{BondarV16}
that $\Gamma_1(\mathcal A)$ is strongly connected if and only if every subset is reachable
by a concatenation of words of defect one. Bondar et al.~\cite{BondarCV23}
gave an algorithm to compute $\Gamma_k(\mathcal A)$ with running time~$O(|\Sigma||Q|^{2k} \log(|Q|))$.

Associating to an edge $(p,q)$ in a given digraph the simple idempotent $\SI{p}{q}$
and investigating the semigroup generated by them has been done in~\cite{CameronRGM2017,ALCO_2019__2_5_711_0}.
In these works, the authors start with a digraph and construct a semigroup, whereas
our digraphs $I(S)$ arise from a given semigroup~$S$.

\IN{$k$-regular set of words}
\IN{strongly transitive automaton}
Rystsov~\cite{Rystsov95a,Rystsov1995a} called a set $W \subseteq \Sigma^*$
of words \emph{$k$-regular} w.r.t.~an automaton $\semiautomat{A}{Q}{\Sigma}{\delta}$  
if it contains the empty word, $W \subseteq \Sigma^{< |Q|}$
and for each $p, q \in Q$ there are exactly $k$ words in $W$
that map $p$ to~$q$. He has shown (using \cref{sync_max:eqn:rt_bound} and other techniques)
that the reset length of a synchronizing automaton admitting a $k$-regular set of words 
for some $k \ge 1$ is bounded by $2(n-1)^2$.
Rystsov~\cite{Rystsov95a,Rystsov1995a} conjectured that every 
strongly connected automaton admits a $k$-regular set of words for some $k \ge 1$.
A similar notion was introduced by Carpi and D'Alessandro~\cite{DBLP:journals/acta/CarpiD09} 
who called an $n$-state automaton $\semiautomat{A}{Q}{\Sigma}{\delta}$ \emph{strongly transitive}
if it admits words $w_1, \ldots, w_n$ such that for every pair of states $p,q \in Q$
we have $\delta(p, w_i) = q$ for some $i \in \{1,\ldots,n\}$. They proved
that if such an automaton is synchronizing, then it has a synchronizing word of length
at most $(n - 2)(n + L - 1) + 1$ where $L = \max\{ |w_1|, \ldots, |w_n| \}$.
This result has been extended in~\cite{DBLP:journals/tcs/CarpiD17}.
Carpi and D'Alessandro wrote the survey~\cite{DBLP:journals/jalc/CarpiD19}
on their findings. They conjectured~\cite{DBLP:journals/acta/CarpiD09}
that every strongly connected synchronizing
automaton admits such a sequence with $L < kn$ for some $k \ge 1$.
Both conjectures have been refuted by Berlinkov~\cite{DBLP:journals/ijfcs/Berlinkov11}.
He gave the modified conjecture that in a strongly connected synchronizing automaton 
there exists an augmentative set whose words have length at most~$2n$.